\section{¿Por qué los modelos de difusión pueden modelar distribuciones complejas}

\subsection{Gaussiano de un paso vs. Gaussiano de múltiples pasos}

El conferenciante planteó en clase una pregunta profunda: dado que las transiciones del proceso hacia adelante son gaussianas y la transición del proceso hacia atrás también se modela como gaussiana, ¿la distribución de verosimilitud completa $p_\theta(x_0|x_T)$ también es una distribución gaussiana? Si lo fuera, los modelos de difusión no estarían tan limitados como los VAE.

La respuesta es \textbf{no}.

\begin{importantbox}{La no linealidad introduce complejidad}
Aunque cada transición hacia atrás individual $p_\theta(x_{t-1}|x_t)$ es una distribución gaussiana, la distribución de verosimilitud completa $p_\theta(x_0|x_T) = \int p_\theta(x_0|x_1) p_\theta(x_1|x_2) \cdots p_\theta(x_{T-1}|x_T) \, dx_{1:T-1}$ \textbf{no} es una distribución gaussiana. La clave está en que la \textbf{media} de cada paso gaussiano se predice mediante una red neuronal (una función no lineal).

A diferencia del proceso hacia adelante —donde la media es una función lineal de $x_{t-1}$ ($\sqrt{1-\beta_t}\, x_{t-1}$), por lo que la combinación de múltiples pasos en el proceso hacia adelante sigue siendo gaussiana—, en el proceso hacia atrás la media es una función no lineal de $x_t$ (a través de la red neuronal). Esto permite que la combinación de múltiples pasos exprese cualquier distribución compleja.
\end{importantbox}

\begin{warningbox}{Proceso hacia adelante lineal vs. proceso hacia atrás no lineal}
Esto es fundamental para entender la capacidad de expresión de los modelos de difusión:
\begin{itemize}
    \item \textbf{Proceso hacia adelante}: La media es una función \textbf{lineal} de la entrada $\rightarrow$ la combinación de múltiples pasos sigue siendo gaussiana $\rightarrow$ se puede derivar una distribución de salto en forma cerrada.
    \item \textbf{Proceso hacia atrás}: La media es una función \textbf{no lineal} de la entrada (red neuronal) $\rightarrow$ la combinación de múltiples pasos \textbf{no} es gaussiana $\rightarrow$ la distribución de verosimilitud puede ser arbitrariamente compleja.
\end{itemize}
Los modelos de difusión encuentran un equilibrio inteligente entre ``paso simple en cada etapa de entrenamiento'' y ``fuerte capacidad de expresión global''.
\end{warningbox}

Esto explica también por qué los modelos de difusión requieren múltiples pasos: cuanto más es el número de pasos, más profundas son las transformaciones no lineales y mayor es la capacidad de expresión. Esto es análogo a la razón por la que las redes neuronales profundas tienen mayor capacidad de expresión que las redes superficiales.

\subsection{Resumen de la sección}

Aunque cada transición del modelo de difusión es gaussiana, debido a que el proceso hacia atrás utiliza una red neuronal no lineal para predecir la media, la distribución de verosimilitud global resultante de la combinación de múltiples pasos puede modelar distribuciones arbitrariamente complejas. Esta es la razón fundamental por la cual los modelos de difusión son superiores a los VAE simples.

%% ========== Novena sección: Conexión entre predicción de ruido y emparejamiento de puntuaciones ==========


